\section{Forward-simulated short-premium scenarios}\label{sec:scenarios}

To exercise the calibrated framework as a synthetic-data generator and stress test, we forward-simulated 1{,}000 Goldman Sachs (GS) price paths over a 31 trading-day horizon and tracked the path-conditional fair value of two real GS option contracts pulled directly from the 2026-04-28 ladder capture: a short put at $K_p = \$890$ (market $\Delta = -0.295$, mid \$16.51, market IV $31.3\%$) and a short call at $K_c = \$970$ (market $\Delta = +0.328$, mid \$16.09, market IV $28.9\%$), both expiring 2026-05-29 with the underlying spot at $S_0 = \$926.38$. Each path drew from the JumpHMM marginal pretrained on GS returns. Conditional on the path, we evolved a Cox-Ingersoll-Ross variance level $v_t$ with mean-reversion $\kappa = 2.0$, vol-of-vol $\sigma_v = 0.5$, and leverage coupling $\rho = -0.6$, and at every trading day along every path we re-priced both contracts via 201-step Leisen-Reimer American with $\sigma_t^2 = v_t \cdot \psi_{\mathrm{NN}}((K/S_t)_{\mathrm{std}}, (T-t)_{\mathrm{std}})$, the per-ticker neural smile/term-structure surface from the ladder calibration. We adopted Leisen-Reimer over the standard CRR lattice because the Peizer-Pratt construction places the strike $K$ at a fixed lattice position by design at every $(S, \sigma)$ configuration, which gave smooth Greeks via finite differences without the staircase aliasing that CRR's recombining tree produced when bumps moved $K$ across node boundaries. The seller received the market mid as the entry premium, and the model marked the position to its Leisen-Reimer fair value at every subsequent step; we assumed no scheduled earnings event during the 31-day window so the only sources of variation were diffusion, jumps, and the variance dynamics themselves. The model's $t=0$ fair value sat essentially at the market mid on both legs (put: \$17.26 vs.\ \$16.51, edge $+\$0.75$; call: \$16.07 vs.\ \$16.09, edge $-\$0.02$; Supplementary Table~\ref{tab:gs_scenario}), so the simulated short was being sold at par with the model's own fair value at the entry capture and the realised path-conditional P\&L could be read directly as a return.

The forward-simulation engine is summarised in Algorithm~\ref{alg:forward_sim}. Each simulated trading day along each path produced a $(S_t, \sigma_t)$ pair from which the LR pricer returned a fair value for the contract; the same $(S_t, \sigma_t)$ pair fed Algorithm~\ref{alg:greeks} for the central finite-difference Greeks. The two algorithms shared the lattice constructor and the per-ticker $\psi_{\mathrm{NN}}$ surface, so the Greek bumps inherited the same shape function used for the base mark and did not require a separate calibration.

\begin{algorithm}[ht]
\caption{Forward simulation of path-conditional option premia}\label{alg:forward_sim}
\begin{algorithmic}[1]
\Require JumpHMM marginal for ticker $i$; per-ticker $\psi_{\mathrm{NN}}^{(i)}$ and level $\theta_i$; contract $(K, T)$; CIR parameters $(\kappa, \sigma_v, \rho)$; spot $S_0$; risk-free rate $r$; number of paths $N_p$; LR tree depth $N_{\mathrm{LR}}$; step $\Delta t = 1/252$.
\Ensure $\{S_t^{(j)}, v_t^{(j)}, \sigma_t^{(j)}, P_t^{(j)}\}$ for $j = 1,\ldots,N_p$ and $t = 0, \ldots, T$.
\State $v_0 \gets \theta_i \cdot \psi_{\mathrm{NN}}^{(i)}\!\big(\widetilde{\ln(K/S_0)},\, \widetilde{\ln T}\big)$ \Comment{equilibrium initialization, eq.~\eqref{eq:v0}}
\For{$j = 1$ \textbf{to} $N_p$}
  \State Draw price path $\{S_t^{(j)}\}_{t=0}^{T}$ from the JumpHMM marginal
  \State $v_0^{(j)} \gets v_0$
  \For{$t = 0$ \textbf{to} $T-1$}
    \State $\sigma_t^{(j)} \gets \sqrt{v_t^{(j)} \cdot \psi_{\mathrm{NN}}^{(i)}\!\big(\widetilde{\ln(K/S_t^{(j)})},\, \widetilde{\ln(T-t)}\big)}$
    \State $P_t^{(j)} \gets \mathrm{LR}_{N_{\mathrm{LR}}}\!\big(S_t^{(j)},\, K,\, \sigma_t^{(j)},\, T-t,\, r\big)$ \Comment{Leisen-Reimer American}
    \State Draw $Z_S^{(j)}$ from the JumpHMM return innovation; draw $Z_\perp \sim \mathcal{N}(0,1)$ independent
    \State $Z_v^{(j)} \gets \rho\, Z_S^{(j)} + \sqrt{1 - \rho^2}\, Z_\perp$ \Comment{leverage coupling}
    \State $v_{t+1}^{(j)} \gets \big| v_t^{(j)} + \kappa(\theta_i - v_t^{(j)})\,\Delta t + \sigma_v \sqrt{\max(v_t^{(j)},\,0)}\,\sqrt{\Delta t}\, Z_v^{(j)} \big|$
  \EndFor
  \State $P_T^{(j)} \gets \max(\phi (S_T^{(j)} - K),\, 0)$ \Comment{$\phi=+1$ call, $\phi=-1$ put}
\EndFor
\State \Return $\{S_t^{(j)}, v_t^{(j)}, \sigma_t^{(j)}, P_t^{(j)}\}$
\end{algorithmic}
\end{algorithm}

\begin{algorithm}[ht]
\caption{Path-conditional Greeks via central finite differences on the LR pricer}\label{alg:greeks}
\begin{algorithmic}[1]
\Require State $(S, \sigma, \tau)$ with $\tau = T - t$; contract $(K, r)$; LR depth $N_{\mathrm{LR}}$; spot bump $h_S$ (relative); IV bump $h_\sigma$ (absolute).
\Ensure $\Delta$, $\Gamma$, Vega.
\State $P_+^S \gets \mathrm{LR}_{N_{\mathrm{LR}}}(S(1+h_S),\, K,\, \sigma,\, \tau,\, r)$
\State $P_-^S \gets \mathrm{LR}_{N_{\mathrm{LR}}}(S(1-h_S),\, K,\, \sigma,\, \tau,\, r)$
\State $P_0  \gets \mathrm{LR}_{N_{\mathrm{LR}}}(S,\, K,\, \sigma,\, \tau,\, r)$
\State $P_+^\sigma \gets \mathrm{LR}_{N_{\mathrm{LR}}}(S,\, K,\, \sigma + h_\sigma,\, \tau,\, r)$
\State $P_-^\sigma \gets \mathrm{LR}_{N_{\mathrm{LR}}}(S,\, K,\, \sigma - h_\sigma,\, \tau,\, r)$
\State $\Delta \gets \dfrac{P_+^S - P_-^S}{2\, h_S\, S}$
\State $\Gamma \gets \dfrac{P_+^S - 2\, P_0 + P_-^S}{(h_S\, S)^2}$
\State $\mathrm{Vega} \gets \dfrac{P_+^\sigma - P_-^\sigma}{2\, h_\sigma}$ \Comment{per unit IV; multiply by $0.01$ for per-$1\%$ IV}
\State \Return $(\Delta,\, \Gamma,\, \mathrm{Vega})$
\end{algorithmic}
\end{algorithm}

We rendered the share-price ensemble alongside the short-put and short-call premium ensembles (Figs.~\ref{fig:gs_put_paths}, \ref{fig:gs_call_paths}). On the worst-5\% paths (red, Fig.~\ref{fig:gs_put_paths}), GS drifted toward and through $K_p$, the short-put premium climbed many-fold above its entry price, and the position closed with a mean P\&L of $-\$113.66$ on those paths. On the top 1--5\% paths (green, Fig.~\ref{fig:gs_call_paths}; the wildest 1\% were excluded for legibility), GS pushed past $K_c$ and the short-call premium spiked despite the calendar pulling extrinsic value toward zero; mean P\&L on the top-5\% paths was $-\$176.75$, $1.55{\times}$ the mean loss on the worst put paths. On both panels the median path stayed inside the 25--75\% band and closed flat at the premium received; the short-premium trade made money on the typical path, with both legs carrying a long left tail.

\begin{figure}[ht]
\centering
\includegraphics[width=\textwidth]{sections/figures/gs/gs_short_put_paths.pdf}
\caption{Forward GS share-price paths (left) and short put premium paths (right) for the real GS 2026-05-29 \$890 contract, 1{,}000 simulated trajectories over 31 days. Red: worst 5\% by terminal price. Black solid: median; shaded band: 25--75\% IQR. Dotted line on the left marks the strike $K_p = \$890$; blue dashed line on the right marks the market mid received at sale, \$16.51.}
\label{fig:gs_put_paths}
\end{figure}

\begin{figure}[ht]
\centering
\includegraphics[width=\textwidth]{sections/figures/gs/gs_short_call_paths.pdf}
\caption{Forward GS share-price paths (left) and short call premium paths (right) for the real GS 2026-05-29 \$970 contract. Green: top 1--5\% by terminal price; the most extreme 1\% are omitted from the overlay so the figure is not dominated by spike outliers, but the underlying summary statistics use the full top 5\%. Dotted line marks the strike $K_c = \$970$; blue dashed line marks the market mid, \$16.09.}
\label{fig:gs_call_paths}
\end{figure}

The short-call leg lost asymmetrically because IV at the contract strike rose as $t \to T$ on every path (Fig.~\ref{fig:gs_iv}). The rise was not a model artefact: it was the empirical short-dated wing-bow that the per-ticker $\psi_{\mathrm{NN}}$ learned from the ladder corpus. As $(T-t)$ shrank at fixed $\log(K/S)$, the standardised moneyness drifted deeper into the wing where $\psi_{\mathrm{NN}}$ was monotonically larger, and a Gaussian diffusion would have underpriced that wing without the inflation, since short-horizon returns were far more leptokurtic than the central-limit-theorem regime of long-horizon returns. The level $v_t$ contributed only a slow mean-reverting drift around $\bar\theta$, so the rise was essentially a smile-shape effect at a fixed point on the surface. On the green tail $S_t$ pushed away from the call strike, dragging the standardised moneyness even further into the wing and amplifying the rise; on the red tail the leverage coupling lifted $v_t$ at the same time as $S_t$ fell, compounding the IV climb.

\begin{figure}[ht]
\centering
\includegraphics[width=\textwidth]{sections/figures/gs/gs_iv_trajectories.pdf}
\caption{Path-conditional implied volatility at the put strike (left) and call strike (right) over the 31-day horizon for the GS scenario. Both panels showed IV rising into expiry: at fixed $\log(K/S)$, shrinking $(T-t)$ pushed the standardised moneyness deeper into the wing of $\psi_{\mathrm{NN}}$, which the per-ticker network learned to be progressively higher and steeper at short horizons. Tail-binned paths (red, green) showed an additional rise driven by $S_t$ moving away from $K$ and, on the put side, leverage coupling lifting $v_t$ on down-paths.}
\label{fig:gs_iv}
\end{figure}

We decomposed the position risk into $\Delta$, $\Gamma$, and Vega along the same paths (Fig.~\ref{fig:gs_greeks}), computed by central finite differences on the same 201-step Leisen-Reimer American pricer used for the base marks (spot bumps of $\pm 1.5\%$ for $\Delta$ and $\Gamma$; $\sigma$ bumps of $\pm 0.005$ for Vega-per-1\%-IV). The long-put $\Delta$ began near $-0.30$ and bifurcated: on the worst-5\% paths it ran to $-1$ as the put pinned ITM at expiry, while on the rest it ran to $0$; the long-call mirror image held at the call strike on the top tail. $\Gamma$ collapsed to zero on every path that drifted away from its strike but spiked sharply on paths that pinned to the strike near expiry, exhibiting the canonical short-gamma blowup geometry that motivates the practitioner aversion to short-premium expiry-day exposure. Vega decayed monotonically toward zero on every path (the $\sqrt{T-t}$ factor dominates), but the worst-5\% put paths and top-5\% call paths held materially more vega than the median through about $t = 20$, which exposed the short-premium trader to vol-of-vol risk on precisely the paths that were already losing on $\Delta$. The signs of the position-level Greeks for the short trader are the negatives of the long Greeks shown.

\begin{figure}[ht]
\centering
\includegraphics[width=\textwidth]{sections/figures/gs/gs_short_greeks.pdf}
\caption{Leisen-Reimer American Greeks ($\Delta$, $\Gamma$, Vega) at the GS put and call strikes over the 31-day horizon, computed by central finite differences on the 201-step Leisen-Reimer pricer at the per-path NN-IV ($h_S = 1.5\% S_t$ for $\Delta$/$\Gamma$, $d\sigma = \pm 0.005$ for Vega). Long-Greek convention: short-position Greeks are the negatives. Top row: $\Delta$. Middle: $\Gamma$, exhibiting the short-gamma spike on paths that pinned to the strike near expiry. Bottom: Vega, decaying toward zero with $\sqrt{T-t}$ on every path. Tail overlays as in Figs.~\ref{fig:gs_put_paths} and~\ref{fig:gs_call_paths}.}
\label{fig:gs_greeks}
\end{figure}

Both legs showed the characteristic short-premium asymmetry (Fig.~\ref{fig:gs_pnl}; moments in Supplementary Table~\ref{tab:gs_scenario}): a thin probability mass piled at the premium ceiling, paired with a long left tail. The short put kept its full \$16.51 premium on $72.2\%$ of paths and closed with a mean P\&L of $+\$1.52$; the short call kept its full \$16.09 premium on $64.3\%$ of paths and its mean P\&L of $-\$9.00$ was net negative, with the worst-case loss reaching $-\$484.57$, $1.50{\times}$ the worst put loss and $30{\times}$ the premium itself. The asymmetry between the two legs reflected the negative skew that GS's per-ticker $\psi_{\mathrm{NN}}$ encodes (higher IV at down-strikes than at equally-far up-strikes) together with the leverage coupling that lifted $v_t$ on down-paths; the put leg was implicitly priced for the down-tail and earned a fair premium for it, while the call leg encountered a fatter empirical right-tail than its strike-symmetric BS-implied premium accounted for.

\begin{figure}[ht]
\centering
\includegraphics[width=\textwidth]{sections/figures/gs/gs_short_pnl_distributions.pdf}
\caption{Terminal P\&L distributions per contract (1{,}000 paths, log-y) for the GS scenario. Both legs showed the short-premium signature: a tall spike of paths keeping the full premium, and a long left tail. The short call's left tail was substantially fatter than the put's despite a smaller premium received.}
\label{fig:gs_pnl}
\end{figure}

\subsection{Cross-ticker validation: LLY}\label{sec:lly_scenario}

To confirm the scenario engine's behaviour was not specific to GS, we re-ran the identical pipeline on a real LLY option pair pulled from the same 2026-04-28 capture: a short put at $K_p = \$825$ (market $\Delta = -0.303$, mid \$23.30, market IV $44.4\%$) and a short call at $K_c = \$940$ (market $\Delta = +0.309$, mid \$20.76, market IV $44.0\%$), both expiring 2026-05-29 with LLY spot $S_0 = \$873.83$. Every other framework choice---the JumpHMM marginal pretrained on LLY returns, the Heston $(\kappa,\sigma_v,\rho) = (2.0,\,0.5,\,-0.6)$ variance dynamics, the 201-step Leisen-Reimer pricer, the per-ticker $\psi_{\mathrm{NN}}$ trained on the same fifteen-date ladder---was held identical to the GS scenario. LLY had dropped about $5\%$ on 04-28 itself, so the contracts were captured against an elevated-IV market state and the per-ticker $\psi_{\mathrm{NN}}$ (which encodes the time-averaged shape across fifteen capture dates) sat below the realised 04-28 IV; the model's $t=0$ fair value therefore came in below the market mid (entry edge $-\$1.35$ put, $-\$1.80$ call; Supplementary Table~\ref{tab:lly_scenario}).

\begin{figure}[ht]
\centering
\includegraphics[width=\textwidth]{sections/figures/lly/lly_short_pnl_distributions.pdf}
\caption{LLY terminal P\&L distributions per contract (1{,}000 paths, log-y), parallel to Fig.~\ref{fig:gs_pnl}. The short-premium asymmetry persisted on a ticker with substantially higher IV regime: short put kept full premium on $77.4\%$ of paths (mean P\&L $+\$14.31$); short call kept full premium on $75.7\%$ of paths (mean P\&L $+\$6.45$) with a worst-case loss of $-\$303.88$, $1.64\times$ the worst put loss and $14.6\times$ the premium itself. Full statistics in Supplementary Table~\ref{tab:lly_scenario}; share-price paths, IV trajectories, and Greek panels in Supplementary §\ref{sec:supp_lly_scenario}.}
\label{fig:lly_pnl}
\end{figure}

The LLY run reproduced the GS result's qualitative pattern (the short call carried fatter tail risk than the short put despite a smaller premium received) on a ticker with materially higher IV regime---LLY at $\sim 44\%$ market IV against GS's $\sim 30\%$---and a different sector with a different jump-state mixture in the JumpHMM marginal. The higher IV regime placed the LLY $30\Delta$ strikes further from spot for the same nominal delta (LLY put $K/S = 0.944$, call $K/S = 1.076$; GS put $K/S = 0.961$, call $K/S = 1.047$), which is why LLY's premium-kept rates ($77\%/76\%$) ran above GS's ($72\%/64\%$): wider strikes were harder to breach within the 31-day window. The cross-ticker robustness of the pattern argued that the asymmetric short-call tail risk reflected the framework's structural negative skew rather than ticker-specific calibration artefacts.

